\subsection{Section summary}\label{subsec:corpus-section-summary}

This section defines the data foundation of the thesis. This thesis uses a corpus of 977 histopathology-focused English papers downloaded from the PMC Open Access dataset. 

Three subsets of papers are derived from this ingested corpus: a 27-PDF subset for document-extraction evaluation (Subsection~\ref{subsec:doc-extraction-eval-set}); a 15-paper knowledge-extraction calibration cluster selected by ILP on mutual entity relatedness (Subsection~\ref{subsec:knowledge-extraction-calibration-set}); and a randomly sampled 15-paper held-out set for the knowledge-extraction generalization test (Subsection~\ref{subsec:knowledge-extraction-eval-set}). The subsets are pairwise disjoint, so that the measurements on one do not affect the others.

This section also clarified the annotation status of the knowledge-extraction reference labels. These labels were generated by an LLM, not expert-annotated, for scalability. These LLM-generated silver labels, therefore, were not treated as ground truth. 

Overall, this section established what data the thesis uses, how the calibration and evaluation subsets were generated from the initial data, and the limitations of the data design. The following sections and chapters build on this foundation: Section~\ref{section:Document_Extraction} describes document extraction, Section~\ref{section:Knowledge_Extraction} describes knowledge extraction from the extracted text from the document-extraction pipeline, Section~\ref{section:Evaluation_Methodology} defines evaluation metrics, Chapter~\ref{chapter:Results} reports the results, and Chapter~\ref{chapter:Discussion} discusses the implications of data limitations on the interpretation of the results.